\medskip

\emph{Les probabilités demandées seront exprimées sous forme de fractions irréductibles}

\medskip

\textbf{Partie A}

\medskip

On lance trois fois de suite une pièce de monnaie bien équilibrée. On note $X$ la variable aléatoire qui compte le nombre de fois, sur les trois lancers, où la pièce est retombée du côté \og Face \fg.

\medskip

\begin{enumerate}
\item Préciser la nature et les paramètres de la loi de probabilité suivie par $X$.



\item Recopier et compléter le tableau suivant donnant la loi de probabilité de $X$
\begin{center}
\begin{tabularx}{\linewidth}{|*{5}{>{\centering \arraybackslash}X|}}\hline

$k$& 0 & 1 & 2& 3 \\ \hline
$P(X = k)$&&&&\\ \hline
\end{tabularx}
\end{center}
\end{enumerate}

\medskip

\textbf{Partie B}

\medskip

Voici les règles d'un jeu où le but est d'obtenir trois pièces du côté \og Face \fg{} en un ou deux essais :

\begin{itemize}[label=$\bullet~$]
\item On lance trois pièces équilibrées :
	\begin{itemize}[label=$\circ~$]
	\item Si les trois pièces sont tombées du côté \og Face \fg, la partie est gagnée ;
	\item Sinon, les pièces tombées du côté \og Face\fg{} sont conservées et on relance
celles tombées du côté \og Pile \fg.
	\end{itemize}
\item La partie est gagnée si on obtient trois pièces du côté \og Face \fg, sinon elle est
perdue.
\end{itemize}

On considère les évènements suivants :
\begin{itemize}[label=$\bullet~$]
\item $G$ : \og la partie est gagnée \fg.

Et pour tout entier $k$ compris entre 0 et 3, les évènements:
\item $A_k$ : \og $k$ pièces sont tombées du côté \og Face\fg{} au premier lancer \fg.
\end{itemize}

\medskip

\begin{enumerate}
\item Démontrer que $P_{A_1}(G) = \dfrac14$.
\item Recopier et compléter l'arbre pondéré ci-dessous:

\begin{center}
\pstree[treemode=R,levelsep=3cm]{\TR{}}
{\pstree{\TR{$A_0~~$}\taput{$\frac18$}}
	{\TR{$G$}
	\TR{$\overline{G}$}
	}
	\pstree{\TR{$A_1~~$}\taput{$\frac38$}}
	{\TR{$G$}
	\TR{$\overline{G}$}
	}
	\pstree{\TR{$A_2~~$}\taput{$\frac38$}}
	{\TR{$G$}
	\TR{$\overline{G}$}
	}
	\pstree{\TR{$A_3~~$}\tbput{$\frac18$}}
	{\TR{$G$}	
	}
}
\end{center}

\item Démontrer que la probabilité $p$ de gagner à ce jeu est $p = \dfrac{27}{64}$
\item La partie a été gagnée. Quelle est la probabilité qu'exactement une pièce soit tombée du côté \og Face \fg{} à la première tentative ?
\item Combien de fois faut-il jouer à ce jeu pour que la probabilité de gagner au moins une partie dépasse $0,95$ ?
\end{enumerate}

\bigskip
